\section{Components of a drip system}
A \kw{drip} (trickle) system delivers water in small, frequent amounts through emitters placed on or near the plants. The parts, in the order of the flow, are:
\begin{exlist}
\item \kw{The pump and the control head.} The pump gives the pressure; the head holds the valves, the pressure regulator, the flow meter and the gauges, and the fertigation unit.
\item \kw{The filters.} They protect the emitters from clogging. A \emph{screen} filter (mesh 120 to 200) removes sand and particles; a \emph{disc} filter has stacked grooved discs and holds organic matter; a \emph{media} (sand) filter holds algae and organic material; a \emph{hydrocyclone} removes sand by centrifugal separation. Water from a well needs a cyclone and a screen; surface water needs media and screen filters.
\item \kw{The mainline and submains.} Pipes of PVC or polyethylene that carry the water to the blocks of the field.
\item \kw{The laterals.} Thin polyethylene pipes (12 to 20 mm) along the rows, which carry the emitters.
\item \kw{The emitters.} Devices that reduce the pressure and give a small, steady flow (1 to 8 L\,h$^{-1}$).
\end{exlist}

\section{Types of emitters}
Emitters are classified by their construction and by the way they respond to the pressure. \emph{Online} emitters are pushed into the lateral; \emph{inline} emitters are built into the pipe (for rows of annual crops); \emph{drip tape} has thin walls and emitters moulded inside (for one season); \emph{micro-sprinklers} and \emph{bubblers} give a larger flow over a larger area (for trees). By the hydraulics, an emitter is \emph{non-compensating} (the flow varies with the pressure) or \emph{pressure-compensating} (a flexible diaphragm keeps the flow nearly constant over a range of pressure, about 50 to 400 kPa).
\par\noindent The discharge follows the \kw{emitter equation}
\begin{equation*}\boxed{q=K_d\,h^{x}}\end{equation*}
where $q$ is the flow (L\,h$^{-1}$), $h$ the pressure head, $K_d$ a constant of the emitter and $x$ the \kw{emitter exponent}: $x=1$ for a laminar flow in a long path (very sensitive to pressure and temperature), $x=0.5$ for a turbulent orifice, $0.3$ to $0.4$ for a vortex or a long-path labyrinth and $0$ to $0.1$ for a compensating emitter. A small $x$ gives a better uniformity on sloping or long laterals.

\section{Manufacturing variation and the choice of an emitter}
No two emitters are exactly alike. The manufacturing \kw{coefficient of variation} $C_v=s/\bar q$ (the standard deviation divided by the mean flow of a sample) is 0.03 for a good emitter and 0.10 or more for a poor one. It enters the emission uniformity (Section~\ref{ch:eu}). The choice depends on the crop (spacing and wetted area), the water quality (a large passage clogs less), the slope, the pressure available and the cost.

\section{Advantages and disadvantages of drip irrigation}
\textbf{Advantages:} a high application efficiency (85 to 95 per cent) because only the root zone is wetted; little evaporation from the soil between the rows; a smaller leaf wetting and less disease; the water can carry fertiliser (\emph{fertigation}) in the amount and at the time that the crop needs; the use of saline water is possible because the salt is moved to the edge of the wetted bulb and the root zone stays wet; it works on irregular land and light soils; a low labour cost when automated.
\par\noindent \textbf{Disadvantages:} a high capital cost; the clogging of the emitters by sand, algae, iron or calcium (the largest practical problem); the need for filtration and maintenance; the accumulation of salt at the edge of the wetted zone, which a rain moves to the roots; the damage to laterals by rodents and machines; and a limited root volume, so that the crop has less reserve if the water stops.

\section{Clogging and its prevention}
Emitters clog by \emph{physical} causes (sand, silt, particles of plastic), \emph{chemical} causes (precipitation of calcium carbonate at pH above 7.5, of iron and manganese oxides, of phosphate fertiliser with hard water) and \emph{biological} causes (algae and bacterial slime, especially in warm water). The remedies are filtration, flushing of the laterals at the end of the day, treatment with chlorine (against the biological films) and with acid (against the chemical deposits), and a choice of fertiliser that does not precipitate.

\section{The shape of the wetted volume}
Water from a point source forms a bulb whose shape depends on the soil. In a sand the water goes down (a narrow, deep bulb) and in a clay it goes sideways (a wide, shallow one); the same discharge and time give a different diameter, so that the spacing of emitters is closer in light soils. The wetted area should cover 30 to 40 per cent of the area of the crop for a row crop, and at least 33 per cent for orchards in a humid climate and 50 per cent in a dry one; the design of Section~\ref{sec:bulb} gives the spacing of emitters that does this.
